\section{실험 방법}\label{sec:methods}
모든 제어기는 같은 측정 온도, \SI{1}{s} 표본 주기, \SI{3600}{s} 구간, \SI{20}{\celsius}/\SI{50}{\celsius}의 초기/기준 온도, 공칭 한계 $[-2000,3000]$ W를 사용한다. 시험 시나리오는 매개변수별 7단계 변화, 세 가지 결합 사례, \SI{1800}{s}에서의 $-\SI{700}{W}$ 계단 부하, $\sigma=\SI{0.15}{\celsius}$ 가우시안 측정 잡음, 제한 한계 $[-500,1200]$ W를 포함한다.

\subsection{동일 예산 튜닝}
제어기마다 64개의 결정론적 저불일치 후보를 평가한다. PID 범위는 $K_p\in[100,1200]$, $K_i\in[0.2,12]$, $K_d\in[0,3000]$이고, LQR은 $Q\in[0.05,20]$, $R\in[10^{-7},2\times10^{-3}]$, LQI는 $Q_x\in[0.05,20]$, $Q_z\in[10^{-5},2\times10^{-2}]$와 같은 $R$ 구간이다. $K_d$ 이외에는 로그 척도를 쓴다. 모든 후보는 공칭 추종, $-\SI{350}{W}$ 부하, 웜스타트 시나리오를 공유한다. 제어기별 훈련 상위 8개 후보를 동일한 네 가지 $\pm15\%$ 검증 사례에 투입한다. 검증에서 한 번 선택하며 보류 시험 결과로 선택을 바꾸지 않는다.

목적함수는 정규화 IAE, 절대 말단 오차의 2배, 포화 비율의 0.15배, 정규화 RMS 제어입력의 0.05배의 합이다. 동일 후보 수는 계산량을 제한하지만 2차원과 3차원 탐색 난이도를 같게 만들지는 않는다.

지표는 $N+1$ 끝점 표본에 대해 $N$개의 유지 구간을 사용한다. 잡음 분석은 시드 $20261116+i$를 갖는 독립 ChaCha8 스트림 100개를 사용하고 시드별로 제어기를 짝지어 비교한다. 신뢰구간은 반복 차이와 $t_{99}=1.984$를 사용한다. 통계적 식별 가능성과 실용적 크기를 분리해 보고한다.
